%% Section 1: Introduction
%% Every quantitative claim annotated with its frozen source artifact.
%% Source authority: artifacts/ directory, git commit 9bc30db

\section{Introduction}
\label{sec:intro}

Large-scale GPU supercomputers routinely operate at the boundary of their
thermal design envelope. Individual compute nodes exhibit measurable and
persistent differences in thermal behavior under identical workloads: a node
that runs measurably warmer than its neighbors during one multi-hour job tends
to do so in subsequent jobs of similar character. This heterogeneity is not
random noise — it reflects durable physical asymmetries in node-level cooling
efficiency, junction characteristics, and local airflow that persist across
months and across calendar years~\cite{TODO:thermal-heterogeneity-hpc}.

Current production schedulers for large-scale HPC systems assign jobs to nodes
based on availability, topology, and performance characteristics, but treat
nodes as thermally interchangeable~\cite{TODO:hpc-scheduling-survey}.
The persistent thermal structure of a fleet is therefore an unexploited signal
for workload placement. If a placement policy routes thermally sensitive
workloads preferentially to nodes with persistent thermal headroom — nodes
whose historical fingerprints indicate consistently cooler operation — it may
delay the accumulation of high-temperature node-hours that otherwise reduce
feasible capacity. Whether this advantage is measurable, reproducible, and
attributable to node-specific identity rather than cabinet-level structure is
the question this paper addresses.

\paragraph{This work.}
We characterize thermal fingerprints longitudinally on the 9{,}856-node
% source: artifacts/phase1/analysis/phase3b3f_results.json (n_all_nodes=9856)
Frontier supercomputer~\cite{TODO:frontier}, using two years of job-level
power and temperature telemetry. For each node, we define a scalar
\emph{thermal fingerprint} from 2024 training-period telemetry. These
fingerprints prospectively predict nodes' mean relative thermal position in
2025 with Pearson $r = 0.926$
% source: artifacts/revision2/m1_sensitivity_summary.json
%   key_stats.individual_F_pearson_r = 0.9258770969716581
and frozen $R^2 = 0.848$,
% source: artifacts/revision2/m1_sensitivity_summary.json
%   key_stats.r2_frozen_individual_F = 0.8482261482495985
establishing that thermal identity is a durable property of individual nodes
across calendar-year timescales. The fingerprint signal is not reducible to
cabinet or chassis location: location-only models predict 2025 thermal
positions with frozen $R^2 = 0.058$, a difference of $0.790$ from the
% source: artifacts/revision2/m1_sensitivity_summary.json
%   key_stats.r2_frozen_location_only = 0.05846979220614623
%   key_stats.difference_in_frozen_predictive_R2 = 0.7897563560434523
individual-fingerprint model; a within-chassis permutation test yields
$p = 1/2001 \approx 0.00050$ against the null of no node-specific signal.
% source: artifacts/revision2/m1_sensitivity_summary.json
%   key_stats.within_chassis_permutation = "p < 1/2001 (~0.0005)"

We then quantify the scheduling-capacity advantage through a
\emph{pre-specified, frozen held-out challenge}: 25 scenarios (sc30--sc54), 8
comparisons in a factorial placebo structure, all executed under a protocol
frozen at a fixed Git commit before any scenario ran. We use a
\emph{constraint-aware operating-envelope replay model}: the full 2{,}637-job
% source: artifacts/phase1/analysis/phase3b_mc_results.json (n_base_jobs=2637)
historical base workload is replayed exactly, and incremental synthetic
workload is added to measure the capacity-yield crossing point — the load
multiplier at which each policy's thermal budget becomes binding — under each
placement strategy.

\paragraph{Key results.}
At the pre-specified $1.25{\times}$ candidate choice ratio, fingerprint-guided
placement (full\_fp) produced a median conservative bracket separation of
$+6.71$~pp versus random placement
% source: artifacts/revision2/challenge_analysis.json
%   FP_vs_Random r=1.25: median_delta_lo = 6.706290619727605
(bootstrap 95\%~CI $[{+}3.92, {+}8.72]$~pp; 22 of 25 held-out scenarios
% source: artifacts/revision2/challenge_analysis.json
%   bootstrap95_delta_lo = [3.9196535626004456, 8.72485557161271]
%   n_definitely_A = 22
definitely fingerprint-first). The conservative bracket separation
$\delta_{\mathrm{lo}} = C_{\mathrm{lower}}(\mathrm{FP}) -
C_{\mathrm{upper}}(\mathrm{Random})$ is a strict lower bound on the capacity
separation: a positive value means the lower bound of FP's capacity bracket
exceeds the upper bound of random placement's bracket.

To test whether this advantage reflects individual node identity or merely
cabinet-level structure, we compare against a within-cabinet permutation
placebo that randomly reassigns fingerprint scores among nodes within each
cabinet, preserving the cabinet-level distribution while breaking individual
node-fingerprint associations. Fingerprint-guided placement beats this placebo
in 20 of 25 scenarios ($+3.54$~pp; CI $[{+}0.56, {+}4.44]$~pp),
% source: artifacts/revision2/challenge_analysis.json
%   FP_vs_WithinCabPlacebo r=1.25: median_delta_lo = 3.537760317644416
%   bootstrap95_delta_lo = [0.5580646330144061, 4.435295391457885]
%   n_definitely_A = 20
providing held-out evidence that individual node identity contributes to
scheduling advantage beyond cabinet-level score composition.

Results are robust to $\pm10\%$ cabinet-envelope calibration error: the sign
of the median conservative bracket separation is positive at all five tested
scale factors.
% source: artifacts/revision2/sensitivity_cabinet_scale.json
%   sign positive at scales [0.90, 0.95, 1.00, 1.05, 1.10]
The qualitative direction is stable across six pre-specified beta estimator
variants spanning $[0.96, 1.17]$ with no sign reversal.
% source: artifacts/revision2/sensitivity_loo_beta.json
%   6 betas: primary (1.114), frozen (1.125), 2025-in-sample (0.96),
%   jackknife (1.083), CF-P5 (1.066), CF-P95 (1.169); no sign reversal

\noindent\textbf{Limitation (retained here):}
Fingerprint-guided placement was \emph{not} statistically separated from the
thermally informed structural-position template at the primary choice ratio
(16/5/4 scenarios; median $+0.74$~pp, 95\%~CI $[{-}0.82, {+}1.50]$~pp,
% source: artifacts/revision2/challenge_analysis.json
%   FP_vs_PositionTemplate r=1.25: n_A=16, n_I=5, n_B=4
%   median_delta_lo = 0.7411474885832706
%   bootstrap95_delta_lo = [-0.8163, +1.5035]
spanning zero). The position template is not a purely topological index; it
reflects average historical fingerprint scores at each physical structural
position. We do not claim fingerprints outperform all thermally informed
placement strategies.

\paragraph{Contributions.}
This paper makes five contributions:

\begin{itemize}

\item \textbf{Large-scale longitudinal fingerprint characterization.}
We characterize persistent node-level thermal heterogeneity across a production
GPU supercomputer: a 2{,}000-node cohort spanning 77~cabinets,
% source: artifacts/phase1/analysis/phase3b3f_results.json (n_all_nodes=9856, n_cab=77)
% 2,000-node cohort: consistent across phase3b artifacts
462 qualifying training jobs, and 2{,}637 test jobs spanning two calendar
% source: artifacts/revision2/loo_beta_results.json (training_data.n_jobs_attempted=462)
% source: artifacts/phase1/analysis/phase3b_mc_results.json (n_base_jobs=2637)
years. Fingerprints prospectively predict 2025 thermal positions from 2024
telemetry ($r = 0.926$, frozen $R^2 = 0.848$; \S\ref{sec:fingerprints}).
% source: artifacts/revision2/m1_sensitivity_summary.json

\item \textbf{Replay with validated cabinet-power reconstruction.}
We design and validate a stateful operating-envelope replay model that
reconstructs cabinet-level thermal budgets from production telemetry and adds
incremental synthetic workload to measure capacity-yield crossings.
Cabinet power estimation mean absolute error is 0.37~kW (0.8\% of the
% source: artifacts/revision1/results/power_validation_60job.json
%   overall.mae_kW = 0.3722548980702776; P99 spread = 48.21 kW; 0.77% ~ 0.8%
inter-cabinet P99 spread; \S\ref{sec:model}).

\item \textbf{Pre-specified, frozen challenge protocol.}
We freeze the scientific runner, all comparison hypotheses, beta estimator,
bootstrap parameters, and interpretation rules at a fixed Git commit before any
held-out scenario runs. A 16/16 preflight gate verifies simulation equivalence,
% source: artifacts/revision2/challenge_preflight.json (n_conditions=16, overall_pass=True)
analysis correctness, and seed integrity before execution; a separate completion
audit verifies output integrity after all scenarios complete and before the
% source: artifacts/revision2/challenge_completion_audit.json (audit_pass=True)
frozen analysis runs. This design separates prespecification from execution and
provides verifiable provenance for all results (\S\ref{sec:protocol}).

\item \textbf{Factorial placebo structure for node-specificity.}
Eight pre-specified comparisons (\S\ref{sec:results}) isolate the contribution
of individual node identity from cabinet-level score structure, arbitrary global
ranking, and spreading mechanics. The pattern of results distinguishes a
node-specific scheduling signal from cabinet-level composition effects within
the replay-model scope.

\item \textbf{Pre-specified FPSpread crossover: choice-set-size dependence.}
Fingerprint spreading (FPSpread) underperforms direct fingerprint selection at
constrained choice ratios ($r = 1.25$--$2.0$) but outperforms at large choice
sets ($r \geq 8.0$, unrestricted). This crossover was pre-specified before
execution and is replicated exactly in the 25 held-out scenarios
(\S\ref{sec:results}).

\end{itemize}

\paragraph{Paper organization.}
\S\ref{sec:background} reviews related work on thermal management in HPC,
workload fingerprinting, and preregistered measurement studies.
\S\ref{sec:fingerprints} characterizes thermal fingerprints longitudinally.
\S\ref{sec:model} describes the constraint-aware replay model and its validation.
\S\ref{sec:protocol} details the pre-specified, frozen challenge protocol.
\S\ref{sec:results} presents primary and secondary results.
\S\ref{sec:robustness} presents robustness and sensitivity analyses.
\S\ref{sec:discussion} discusses limitations, scope, and open questions.
\S\ref{sec:conclusion} concludes.
All quantitative claims in this paper are grounded in frozen artifacts; each
quantitative statement in the source carries a \texttt{\% source:} comment
identifying its frozen JSON artifact and field.

\paragraph{Reproducibility.}
The repository and analysis scripts are archived at
\url{https://github.com/K2alyan/thermalshift}
(DOI:~\href{https://doi.org/10.5281/zenodo.23020901}{10.5281/zenodo.23020901}).
Bootstrap seeds and pre-specified analysis parameters are recorded in the frozen
plans and result artifacts, whose integrity and provenance are recorded in
\texttt{final\_output\_audit.json}.
The repository includes derived artifacts sufficient to reproduce the reported
tables, figures, and downstream challenge analyses; reconstruction of
raw-telemetry-dependent measurement stages uses the publicly available OLCF
dataset~\cite{olcf-opendata}
(DOI:~\href{https://doi.org/10.13139/OLCF/3013979}{10.13139/OLCF/3013979}).
